\documentclass[acmsmall,screen,review]{acmart}

\usepackage{proof}
\usepackage{fancyvrb}
\usepackage{minted}
\setminted{fontsize=\small} % Remove this for larger code snippets

\usepackage{latexsym}
\usepackage{tikz}
\usepackage{graphicx}
\usepackage{mathpartir}
\usepackage{color}
\usepackage{pgflibraryarrows}
\usepackage{pgflibraryshapes}
\usepackage{pgfbaseimage}
\usepackage{pgfplots}

\usetikzlibrary{decorations.text}
\usetikzlibrary{trees}
\usetikzlibrary{fadings}

\usepackage{xcolor}

\usepackage{amsthm}

\let\Bbbk\relax
\usepackage{amssymb}

\let\Bbbk\relax
\usepackage{amssymb}

\usepackage[ruled,vlined]{algorithm2e}
\SetAlgoLined
\SetAlgoNoEnd
\SetArgSty{textup}
\SetKwProg{Program}{Program}{\hspace{0.3em}:}{}
\SetKwInOut{Input}{input}\SetKwInOut{Output}{output}
\SetKwIF{Match}{ElseMatch}{Otherwise}{match}{with}{else}{else}{endmatch}

\usepackage{fancyhdr}

\usepackage{booktabs}

\usepackage{multicol}

\setcopyright{acmlicensed}
\copyrightyear{2026}
\acmYear{2026}
\acmDOI{XXXXXXX.XXXXXXX}

\acmJournal{JACM}
\acmVolume{37}
\acmNumber{4}
\acmArticle{111}
\acmMonth{8}

%%
%% Submission ID.
%% Use this when submitting an article to a sponsored event. You'll
%% receive a unique submission ID from the organizers
%% of the event, and this ID should be used as the parameter to this command.
%%\acmSubmissionID{123-A56-BU3}


%%
%% The majority of ACM publications use numbered citations and
%% references.  The command \citestyle{authoryear} switches to the
%% "author year" style.
%%
%% If you are preparing content for an event
%% sponsored by ACM SIGGRAPH, you must use the "author year" style of
%% citations and references.
%% Uncommenting
%% the next command will enable that style.
%%\citestyle{acmauthoryear}


\newcommand{\abs}[4]{{#1}\, #2\! : \! #3.\, #4}
\newcommand{\absu}[3]{{#1}\, #2.\, #3}
\mathchardef\mhyph="2D % Define a "math hyphen"

% unicode
\usepackage[utf8]{inputenc}
\DeclareUnicodeCharacter{3BB}{\ensuremath{\lambda}}
\DeclareUnicodeCharacter{3B4}{\ensuremath{\delta}}
\DeclareUnicodeCharacter{3C9}{\ensuremath{\omega}}
\DeclareUnicodeCharacter{3B3}{\ensuremath{\gamma}}
\DeclareUnicodeCharacter{3C4}{\ensuremath{\tau}}
\DeclareUnicodeCharacter{21D2}{\ensuremath{\To}}
\DeclareUnicodeCharacter{22B2}{\ensuremath{\lhd}}
\DeclareUnicodeCharacter{2080}{\ensuremath{_0}}
%\definecolor{darkred}{RGB}{100,10,10}


\newcommand{\BD}[1]{\textcolor{olive}{[BD: #1]}}
\newcommand{\PA}[1]{\textcolor{cyan}{[PA: #1]}}

\newcommand{\utp}[0]{\mathcal{U}}
\newcommand{\lin}[0]{\lambda^{\iota\nu}}
\newcommand{\interp}[1]{\llbracket #1 \rrbracket}
\newcommand{\interpl}[0]{\llbracket}
\newcommand{\interpr}[0]{\rrbracket}
\newcommand{\choice}[0]{\xi}
\newcommand{\simto}[0]{\triangleright}
\newcommand{\elcap}[0]{\cap}

\newcommand{\dedu}[0]{\triangleright}
\newcommand{\ctxtup}[1]{\ulcorner #1 \urcorner}
\newcommand{\ctxtdn}[1]{\llcorner #1 \lrcorner}
\newcommand{\ctxtmv}[1]{\langle #1 \rangle}

\newcommand{\case}[1]{\noindent\underline{Case #1:}}
\newcommand{\caseb}[0]{\noindent\underline{Case:}}

\newcommand{\pcase}[1]{\noindent\underline{#1:}}


\newcommand{\adbmal}[0]{\reflectbox{$\lambda$}}

\newcommand{\casetri}[0]{\RHD}
\newcommand{\lke}[0]{\lhd}

\newcommand{\lam}[2]{\lambda\, #1.\, #2}
\newcommand{\mal}[2]{\abdmal\, #1.\, #2}
\newcommand{\all}[2]{\forall\, #1.\, #2}
\newcommand{\Lam}[2]{\Lambda\, #1.\, #2}
\newcommand{\thereis}[2]{\exists\, #1.\, #2}
\newcommand{\nuu}[2]{\nu\, #1.\,#2}
\newcommand{\inn}[2]{\langle #1 \rangle_{#2}}
\newcommand{\exel}[2]{\adbmal\,#1.\,#2}
\newcommand{\recc}[2]{\textit{rec}\, #1.\, #2}
\newcommand{\mew}[2]{\mu_{#1}\, #2}
\newcommand{\convl}[0]{\blacktriangleleft}
\newcommand{\convr}[0]{\blacktriangleright}
\newcommand{\conjin}[0]{\smalltriangleleft}
\newcommand{\conjout}[0]{\smalltriangleright}
\newcommand{\dcdot}[0]{\mathbin{\ast}}
\newcommand{\To}[0]{\Rightarrow}
\newcommand{\releq}[0]{\topdoteq}

\newcommand{\omga}[4]{\omega\ #1(#2)\,:\,#3.\ #4}
\newcommand{\clause}[2]{#1 \to #2}

\newcommand{\hysep}[0]{\mathbin{\blacklozenge}}

\newcommand{\todo}[1]{\textbf{[#1]}}
\newcommand{\rref}[1]{\{\ref{#1}\}}

\newtheorem{theorem}{Theorem}
\newtheorem{lemma}[theorem]{Lemma}
\newtheorem{corollary}[theorem]{Corollary}
\newtheorem{definition}[theorem]{Definition}
\newtheorem{notation}[theorem]{Notation}
\newtheorem{proposition}[theorem]{Proposition}


%\newtheorem{definition}[theorem]{Definition}


% \newenvironment{mycomment}
%     {\color{blue} \it
%     }
%     {
%     }
\newcommand{\shift}[2]{{\uparrow^{#1}_{#2}}}
\newcommand{\textblue}[1]{\textcolor{blue}{#1}}

\newcommand{\tl}[0]{\mathsf{L}}
\newcommand{\tn}[0]{\mathsf{N}}
\newcommand{\subt}[4]{(#1)\ #2 <:^{#3} #4}

%%
%% end of the preamble, start of the body of the document source.
\begin{document}

\title{Inferring Coercions for Monadic Programming}

%% \author{Henry Blanchette}
%% \email{blancheh@umd.edu}
%% \affiliation{%
%%   \institution{The University of Maryland}
%%   \city{College Park}
%%   \state{Maryland}
%%   \country{USA}
%% }

%% \author{Ryan Doenges}
%% \email{me@ryandoeng.es}
%% \affiliation{%
%%   \institution{Boston College}
%%   \city{Boston}
%%   \state{Massachusetts}
%%   \country{USA}
%% }

%% \author{Sage Binder}
%% \email{sage.binder@proton.me}
%% \affiliation{%
%%   \institution{The University of Iowa}
%%   \city{Iowa City}
%%   \state{Iowa}
%%   \country{USA}
%% }

%% \author{Robert Zhang}
%% \email{robertz@cs.utexas.edu}
%% \affiliation{%
%%   \institution{The University of Texas at Austin}
%%   \city{Austin}
%%   \state{Texas}
%%   \country{USA}
%% }

%% \author{Caden Parajuli}
%% \email{parajulc@bc.edu}
%% \affiliation{%
%%   \institution{Boston College}
%%   \city{Boston}
%%   \state{Massachusetts}
%%   \country{USA}
%% }

%% \author{Shameer Ahmad}
%% \email{ahmadsk@bc.edu}
%% \affiliation{%
%%   \institution{Boston College}
%%   \city{Boston}
%%   \state{Massachusetts}
%%   \country{USA}
%% }

%% \author{Bill Vu}
%% \email{vubi@bc.edu}
%% \affiliation{%
%%   \institution{Boston College}
%%   \city{Boston}
%%   \state{Massachusetts}
%%   \country{USA}
%% }

%% \author{Aaron Stump}
%% \email{stumpaa@bc.edu}
%% \affiliation{%
%%   \institution{Boston College}
%%   \city{Boston}
%%   \state{Massachusetts}
%%   \country{USA}
%% }


%%
%% By default, the full list of authors will be used in the page
%% headers. Often, this list is too long, and will overlap
%% other information printed in the page headers. This command allows
%% the author to define a more concise list
%% of authors' names for this purpose.
\renewcommand{\shortauthors}{\ }%Blanchette et al.}

%%
%% The abstract is a short summary of the work to be presented in the
%% article.
\begin{abstract}
  This paper shows how to use subtyping to infer uses of the basic
  monad operations.
\end{abstract}

%%
%% The code below is generated by the tool at http://dl.acm.org/ccs.cfm.
%% Please copy and paste the code instead of the example below.
%%
\begin{CCSXML}
<ccs2012>
<concept>
<concept_id>10011007.10011006.10011008.10011009.10011012</concept_id>
<concept_desc>Software and its engineering~Functional languages</concept_desc>
<concept_significance>500</concept_significance>
</concept>
<concept>
<concept_id>10011007.10011006.10011008.10011024.10011033</concept_id>
<concept_desc>Software and its engineering~Recursion</concept_desc>
<concept_significance>500</concept_significance>
</concept>
</ccs2012>
\end{CCSXML}

\ccsdesc[500]{Software and its engineering~Functional languages}
\ccsdesc[300]{Software and its engineering~Recursion}

\keywords{strong functional programming, type-based termination, Mendler-style recursion}

%\received{20 February 2007}
%\received[revised]{12 March 2009}
%\received[accepted]{5 June 2009}

\theoremstyle{definition}
\newtheorem{specification}[section]{Specification}

%%
%% This command processes the author and affiliation and title
%% information and builds the first part of the formatted document.
\maketitle

\section{Introduction}

The basic subtyping axioms for applicatives are:
  \[
\begin{array}{lll}
  \infer[\textit{seq}]{\tl\ (T_1 \to T_2) <: \tl\ T_1 \to \tl T_2}{\ }
  &\ &
  \infer[\textit{pure}]{T <: \tl\ T}{\ }
  \end{array}
  \]

\cite{peytonjones2001}


\section{A Decision Procedure}

\begin{figure}
  \[
  \begin{array}{llll}
    \textit{Types} & S, T & ::= & \tn\ |\ \tl\ T\ |\ T \to T' \\
    \textit{Modes} & \theta & ::= & a\ |\ \tau
  \end{array}
  \]
  \caption{Syntax of types}
  \label{fig:syn}
\end{figure}

Consider this example subtyping:
\[
\tn \to \tn \ \ <: \ \ \tl\ \tn \to \tl\ (\tl\ \tn)
\]
This is valid for applicatives, because we may first use \texttt{fmap}
(which follows from \texttt{pure} and \texttt{seq}) to add $\tl$ to the
domain and codomain parts of the type, and then we may use function subtyping
to expose the codomain part, and add an extra $\tl$ with \texttt{pure}. So the
coercion of $f : \tn \to \tn$ would be
\[
\textit{pure} \cdot (\textit{fmap}\ f)
\]
\noindent Note that with our current set of coercions, we cannot apply the function
subtyping rule at the beginning, because we do not have (contravariantly) $\tl\ \tn <: \tn$.
The challenge of admissibility of transitivity is to derive
coercions that do not use composition.  So we need some form of function subtyping
rule that takes occurrences of $\tl$ in the domain part of the upper bound type ($T_2$ in $T_1 <: T_2$)
into account.  To do this, we adopt a new function subtyping rule that keeps track of how many applications
of the functor would need to be inserted.  

  We work with types $T$ and $T'$ in the syntax shown in
  Figure~\ref{fig:syn}.  We use $\tn$ as a base type (it could be
  the natural numbers), and $\tl$ as an example covariant functor
  (like the list datatype).  The figure also gives the syntax for
  \emph{modes} $\theta$, which play an important role in our decision
  procedure for subtyping.  The mode $\tau$ means a top-level position
  in the types, while $a$ means any position. Figure~\ref{fig:decproc}
  presents that procedure in the form of a proof system for deriving
  judgments of the form $\subt{T}{\theta}{T'}{k}$, where
  $k\in\mathbb{N}$.  An example derivation is shown in Figure~\ref{fig:ex}.

  To handle monads, we would replace the \textit{seq} axiom with
  \[
  \infer[\textit{bind}]{T_1 \to \tl T_2 <: \tl\ T_1 \to \tl T_2}{\ }
  \]
  \noindent Then I am wondering if we could just replace the \textit{fun} rule of Figure~\ref{fig:decproc} with this alternative:

  \[
      \infer[\textit{fun'}]{\subt{T_1\to T_2}{a}{T_1'\to T_2'}{k}}{\subt{T_1'}{\tau}{T_1}{k} & \subt{T_2}{a}{T_2'}{k'} & k' \ge 1 }
      \]

\begin{figure}
  \[
  \begin{array}{lll}
    \infer[\textit{refl}]{\subt{\tn}{a}{\tn}{0}}{\ }
    &\ &
    \infer[\textit{fun}]{\subt{T_1\to T_2}{a}{T_1'\to T_2'}{k}}{\subt{T_1'}{\tau}{T_1}{k} & \subt{T_2}{a}{T_2'}{k'} & k' \ge k }
    \\ \\
    \infer[\textit{seq-dom}]{\subt{\tl\ T}{\tau}{T'}{k+1} }{\subt{T}{\tau}{T'}{k} }
    &\ &
    \infer[\textit{reset}]{\subt{T}{\tau}{T'}{0}}{\subt{T}{a}{T'}{k}}
    \\ \\
    \infer[\textit{seq-cod}]{\subt{T}{a}{\tl\ T'}{k+1} }{\subt{T}{a}{T'}{k} }
    &\ &
    \infer[\textit{cov}]{\subt{\tl\ T}{\theta}{\tl\ T'}{k}}{\subt{T}{\theta}{T'}{k}}
  \end{array}
  \]
\caption{Proposed rules for applicatives}
\label{fig:decproc}
\end{figure}


\begin{figure}
  \[
  \infer{\subt{(\tl\ \tn \to \tl\ \tn)\to \tl\ \tn}{a}{\tl\ (\tn\to\tn) \to \tl (\tl \tn)}{1}}
        {\infer{\subt{\tl\ (\tn\to\tn)}{\tau}{\tl\ \tn\to\tl\ \tn}{1}}
          {\infer{\subt{\tn\to\tn}{\tau}{\tl\ \tn\to\tl\ \tn}{0}}
            {\infer{\subt{\tn\to\tn}{a}{\tl\ \tn\to\tl\ \tn}{1}}
              {\infer{\subt{\tl\ \tn}{\tau}{\tn}{1}}
                     {\infer{\subt{\tn}{\tau}{\tn}{0}}{\infer{\subt{\tn}{a}{\tn}{0}}{\ }}}
                     &\infer{\subt{\tn}{a}{\tl\ \tn}{1}}
                            {\infer{\subt{\tn}{a}{\tn}{0}}{\ }}}}
          }
          & \infer{\subt{\tl\ \tn}{a}{\tl\ (\tl\ \tn)}{1}}
                  {\infer{\subt{\tn}{a}{\tl\ \tn}{1}}
                            {\infer{\subt{\tn}{a}{\tn}{0}}{\ }}}}
        \]
  \caption{Example derivation}
  \label{fig:ex}
\end{figure}



\bibliographystyle{ACM-Reference-Format}
\bibliography{paper.bib}


\end{document}
\endinput
%%
%% End of file `sample-acmsmall-submission.tex'.
